\documentclass[11pt,a4paper]{article}
\usepackage[utf8]{inputenc}
\usepackage[T1]{fontenc}
\usepackage[brazil]{babel}
\usepackage[margin=2.5cm]{geometry}
\usepackage{mathptmx}
\usepackage{enumitem}
\usepackage{booktabs}
\usepackage{tabularx}
\usepackage{array}
\usepackage{hyperref}
\hypersetup{colorlinks=false, pdfborder={0 0 0}}
\usepackage{fancyhdr}

\pagestyle{fancy}
\fancyhf{}
\fancyhead[L]{\small TrackFleet --- Termo de Abertura}
\fancyhead[R]{\small\thepage}
\renewcommand{\headrulewidth}{0.4pt}

\newcolumntype{L}{>{\raggedright\arraybackslash}X}

\begin{document}

\begin{center}
{\Large\bfseries DOCUMENTO DE GOVERNANÇA DO PROJETO}\\[2pt]
{\large\bfseries TrackFleet --- Gestão de Rastreador de Automóveis}\\[4pt]
Disciplina de Gerenciamento de Projetos $\cdot$ Framework PMBOK\textregistered\ 8\textsuperscript{a} Edição $\cdot$ Domínio: Governança\\[2pt]
Integrantes: Enzo Allebrand e Kerlon Ribeiro (dois GPs) $\cdot$ 4\textsuperscript{o} semestre $\cdot$ Data: 13/08
\end{center}

\vspace{2mm}
\hrule
\vspace{4mm}

\noindent\textbf{\large Termo de Abertura}

\vspace{1mm}
\noindent O documento que dá a largada ao projeto: autoriza formalmente o início, dá aos gerentes autoridade para alocar recursos e liga o projeto ao problema que ele resolve.

\section{Campos do Termo de Abertura}
\begin{itemize}[itemsep=4pt]
  \item \textbf{Nome do projeto:} TrackFleet --- Gestão de Rastreador de Automóveis.
  \item \textbf{Problema / oportunidade:} falta de visibilidade em tempo real da frota e de dados objetivos para decisões de manutenção e uso dos veículos, agravada pelo custo elevado das soluções corporativas existentes.
  \item \textbf{Objetivo:} lançar uma plataforma de rastreamento e gestão de frotas acessível para pequenas transportadoras até o fim do semestre letivo.
  \item \textbf{Valor esperado:} redução de custos operacionais e de manutenção corretiva, melhora na segurança dos veículos e dos motoristas, e decisões de gestão de frota baseadas em dados reais.
  \item \textbf{Patrocinador:} Diretoria da transportadora parceira (papel definido pela dupla para fins do exercício).
  \item \textbf{Restrições:} prazo do semestre letivo; orçamento zero para hardware próprio (uso de dispositivos de terceiros via API); conformidade com a LGPD.
  \item \textbf{Premissas:} as transportadoras-alvo já possuem ou têm acesso a rastreadores GPS compatíveis, e os motoristas aceitam o monitoramento dos veículos da empresa.
  \item \textbf{Gerentes do Projeto (GPs):} Enzo Allebrand e Kerlon Ribeiro, em co-gestão.
\end{itemize}

\end{document}
